\subsection{Convolutions}
\lecdate{(Midterm last thurs) Lec 10 - Feb 10 (Week 6)}
we say a fn is cptly supported if its supp is cpt.
\begin{defn}
	if $f,g:\bR\gto\bR$ cptly supp and riemann int'bl, then:
	\begin{equation*}
		(f*g)(x) \ = \ \int_{\bR}f(x-t)g(t)\d t
	\end{equation*}
	is well defn'd, and is called the \textbf{convolution} of $f$ and $g$.
\end{defn} \

\begin{defn}
	a seq $J_{n}:\bR\gto\bR$  of cptly supp riemann int'bl fn's are
	called an \textbf{approximate identity} if:
	\begin{enumerate}[(i)]
		\item $\displaystyle\int_{\bR}I_{n}(x)\d x=1$
		\item $\sup_{n}\displaystyle\int_{\bR}\abs{I_{n}(x)}\d x<\infty$
		\item $\displaystyle\lim_{n\sto\infty}\int_{\abs{x}>\delta}\abs{I_{n}(x)}\d x=0$
			for all $\delta>0$
	\end{enumerate}
\end{defn} \

\begin{thm}
	let $f:[a,b]\gto\bR$ be cts, $f(a)=f(b)=0$. define $\tilde{f}(x)$ as:
	\begin{equation*}
		\tilde{f}(x) \ = \
		\begin{cases}
			f(x) & x\in[a,b] \\ 0 & \trm{otw}
		\end{cases}
	\end{equation*}
	let $I_{n}$ be an approx id.
	then $f_{n}=\tilde{f}*I_{n}$ cvgs uniformly to $\tilde{f}$.
\end{thm} \

\begin{pf}[source=Primary Source Material]
	let $\ep>0$. easy to see: $\tilde{f}$ uni cts, so $\ex\delta>0$ s.t.
	\begin{equation*}
		\abs{\tilde{f}(x)-\tilde{f}(y)}<\ep \qquad \fa\abs{x-y}<\delta
	\end{equation*}
	then:
	\begin{align*}
		f_{n}(x)-\tilde{f}(x)
		& = \int\tilde{f}(x-y)I_{n}(y)\d y - \int\tilde{f}(x)I_{n}(y)\d y \\
		& = \int_{\bR}(\tilde{f}(x-y)-\tilde{f}(x))I_{n}(y)\d y \\
		& = \int_{\abs{y}<\delta}(\tilde{f}(x-y)-\tilde{f}(x))I_{n}(y)\d y
		+ \int_{\abs{y}>\delta}(\tilde{f}(x-y)-\tilde{f}(x))I_{n}(y)\d y
	\end{align*}
	first term:
	\begin{equation*}
		\leq \int_{\abs{y}<\delta}\abs{\tilde{f}(x-y)-\tilde{f}(x)}\abs{I_{n}(y)}\d y
		\leq \ep\cdot\sup_{n}\int_{\bR}\abs{I_{n}(y)}\d y
	\end{equation*}
	2nd term:
	\begin{equation*}
		\leq 2\norm{\tilde{f}}_{\infty}\int_{\abs{y}>\delta}\abs{I_{n}(y)}\d y
	\end{equation*}
	goes to 0 as $n\sto\infty$.
	thus, $\limsup_{n\sto\infty}\sup_{x\in\bR}\abs{f_{n}(x)-\tilde{f}(x)}\leq C\ep$.
\end{pf}

\newpage
\begin{thm}[title=Weierstrass Approximation Theorem]
	sps $f:[a,b]\gto\bR$ cts.
	then $\ex$ seq of polyn's $p_{n}$ s.t. $p_{n}\rightrightarrows f$ for $x\in[a,b]$.
\end{thm} \

\begin{pf}[source=Primary Source Material]
	by subtracting linear fn, wlog $f(a)=f(b)=0$.
	by rescaling, wlog $a=0$ and $b=1$.

	extend $f$ to $\bR$ with:
	\begin{equation*}
		\tilde{f}(x) \ = \
		\begin{cases}
			f(x) & x\in[0,1] \\ 0 & \trm{otw}
		\end{cases}
	\end{equation*}
	define $\beta_{n}(t):=b_{n}(1-t^{2})^{n}\bb{I}_{\abs{t}\leq1}$,
	where $\bb{I}$ indicator fn, and:
	\begin{equation*}
		b_{n} \ = \ \frac{1}{\int_{-1}^{1}(1-t^{2})^{n}\d t}
	\end{equation*}
	we claim the following:
	\begin{enumerate}[(1),itemsep=0pt]
		\item $\beta_{n}$ an approx id
		\item $f_{n}(x)=(\tilde{f}*\beta_{n})(x)$ a polyn $\fa n$ for $x\in[0,1]$
	\end{enumerate}
	together, these prove the statement.

	first, we prove claim 2. we have:
	\begin{equation*}
		f_{n}(x) \ = \ \int_{-1}^{1}\tilde{f}(x-t)\beta_{n}(t)\d t
	\end{equation*}
	by change of variables $u=x-t$, this is equal to:
	\begin{equation*}
		\int_{x-1}^{x+1}\tilde{f}(u)\beta_{n}(x-u)\d u
		= \int_{0}^{1}f(u)\beta_{n}(x-u)\d u
	\end{equation*}
	if $x\in[0,1]$ and $u\in[0,1]$, then $x-u\in[-1,1]$, so this is equal to:
	\begin{equation*}
		b_{n}\int_{0}^{1}f(u)(1-(x-u)^{2})^{n}\d u
	\end{equation*}
	this is a polyn in $x$, since this is equal to:
	\begin{equation*}
		\sum_{k,j}x^{j}\int_{0}^{1}f(u)u^{k}\d u
	\end{equation*}

	now, we prove claim 1. by defn:
	\begin{equation*}
		\int_{\bR}\beta_{n}(t)\d t
		\ = \ \int_{\bR}\abs{\beta_{n}(t)}\d t \ = \ 1
	\end{equation*}
	now, we show:
	\begin{equation*}
		\int_{\abs{t}>\delta}\beta_{n}(t)\d t \ \sto \ 0 \trm{ as } n\sto\infty
	\end{equation*}
	we find an upper bound for $b_{n}$:
	\begin{align*}
		\int_{-1}^{1}(1-t^{2})^{n}\d t
		& \geq \ \int_{-1/\sqrt{n}}^{1/\sqrt{n}}(1-t^{2})^{n}\d t \\
		& \geq \ \int_{-1/\sqrt{n}}^{1/\sqrt{n}}\left(1-\left(\frac{1}{\sqrt{n}}\right)^{2}\right)^{n}\d t \\
		& = \ \frac{2}{\sqrt{n}}\left(1-\frac{1}{n}\right)^{n} \\
		& \geq \ \frac{C}{\sqrt{n}}
	\end{align*}
	holds for all $n\geq2$. thus $b_{n}\leq C\sqrt{n}$, so:
	\begin{align*}
		\int_{\abs{t}>\delta}\beta_{n}(t)\d t
		& \ \leq \ C\sqrt{n}\int_{\abs{t}>\delta}(1-t^{2})^{n}\bb{I}_{\abs{t}\leq1} \\
		& \ \leq \ C\sqrt{n}(1-\delta^{2})^{n}
	\end{align*}
	this goes to 0 as $n\sto\infty$ as needed.

	we have proved both claims, thus the statement is proved.
\end{pf}

why are we defining algebras now lmfao

\begin{defn}
	we say $A\subseteq C_{0}(M)$ \textbf{separates pts} if $\fa p,q\in M$, $\ex f\in A$ s.t.
	$f(p)\neq f(q)$.

	we say $A$ \textbf{vanishes} at $x\in M$ if $f(x)=0$ $\fa f\in A$.
\end{defn} \

\begin{thm}[title=Stone-Weierstrass]
	if $M$ cpt, $A\subseteq C_{0}(M)$ an alg of fn's which separates pts and
	vanishes nowhere, then $\bar{A}=C_{0}(M)$.
\end{thm}

\lecdate{Lec 11 - Feb 12 (Week 6)}

before we prove stone-weierstrass, we need the following lemma.
\begin{lm}
	if $A\subseteq C_{0}(M)$ an alg, then so is $\bar{A}$.
\end{lm}

\newpage
\begin{pf}[source=Primary Source Material]
	we'll just prove $\bar{A}$ closed under products. let $f,g\in\bar{A}$.

	then $\ex f_{n},g_{n}\in A$ s.t.
	$f_{n}\rightrightarrows f$ and $g_{n}\rightrightarrows g$.
	then $f_{n}g_{n}\rightrightarrows fg$, so $fg\in\bar{A}$.
\end{pf} \

\begin{lm}
	if $A\subseteq C_{0}(M)$ an alg which separates pts + vanishes nowhere,
	then $\fa p_{1}\neq p_{2}\in M$ and $c_{1},c_{2}\in\bR$,
	$\ex H\in A$ s.t. $H(p_{1})=c_{1}$ and $H(p_{2})=c_{2}$.
\end{lm} \

\begin{pf}[source=Primary Source Material]
	since $A$ vanishes nowhere, $\ex g_{1},g_{2}$ s.t. $g_{1}(p_{1})\neq0$
	and $g_{2}(p_{2})\neq0$. then $g=g_{1}^{2}+g_{2}^{2}$ doesn't vanish at
	either $p_{i}$.

	let $f\in A$ s.t. $f(p_{1})\neq f(p_{2})$. consider:
	\begin{equation*}
		M =
		\begin{pmatrix}
			g(p_{1}) & f(p_{1})g(p_{1}) \\
			g(p_{2}) & f(p_{2})g(p_{2})
		\end{pmatrix} \qquad
		\det(M) = g(p_{1})g(p_{2})(f(p_{2})-f(p_{1})) \neq 0
	\end{equation*}
	therefore, the eq'n:
	\begin{equation*}
		M
		\begin{pmatrix}
			\xi \\ \eta
		\end{pmatrix}
		\ = \
		\begin{pmatrix}
			c_{1} \\ c_{2}
		\end{pmatrix}
	\end{equation*}
	has a soln $(\xi,\eta)$.
	thus, $H=\xi g+\eta fg\in A$ satisfies $H(p_{1})=c_{1},H(p_{2})=c_{2}$.
\end{pf}

we now prove stone-weierstrass. dont feel like boxing it

first, notice if $f\in\bar{A}$, then $\abs{f}\in\bar{A}$.
to see this, suffices to find $f_{n}\in\bar{A}$ s.t. $f_{n}\rightrightarrows\abs{f}$.
consider:
\begin{equation*}
	I \ = \ [-\norm{f}_{\infty},\norm{f}_{\infty}]
\end{equation*}
by weierstrass approx. thm, $\ex p_{n}(x)$ polys s.t. $p_{n}\rightrightarrows\abs{x}$ on $I$.
consider $f_{n}=p_{n}(f)$. since $\bar{A}$ an alg, $f_{n}\in\bar{A}$ for all $n$.
on the other hand, since $p_{n}\rightrightarrows\abs{x}$ on $I$, then
$f_{n}\rightrightarrows\abs{f}$ on $M$.

next, if $f,g\in\bar{A}$, then:
\begin{equation*}
	\max(f,g) \ = \ \frac{f+g}{2}+\frac{\abs{f-g}}{2} \ \in \ \bar{A} \qquad
	\min(f,g) \ = \ \frac{f+g}{2}-\frac{\abs{f-g}}{2} \ \in \ \bar{A}
\end{equation*}
finally, max/min of finitely many elems of $\bar{A}$ are also in $\bar{A}$.
we now prove the statement.

given $F\in C_{0}(M)$ and $\ep>0$, want to find $G\in\bar{A}$ s.t. $\fa x\in M$:
\begin{equation*}
	F(x)-\ep \ \leq \ G(x) \ \leq \ F(x)+\ep
\end{equation*}
by lemma, $\fa p,q\in M$, $\ex H_{pq}\in\bar{A}$ s.t.:
\begin{equation*}
	H_{pq}(p)=F(p) \qquad H_{pq}(q)=F(q)
\end{equation*}
fix $p\in M$. then $\fa q\in M$, $\ex U_{q}\ni q$ nbhd s.t. $\fa x\in U_{q}$:
\begin{equation*}
	H_{pq}(x) \ > \ F(x)-\ep
\end{equation*}
now, $\set{U_{q}}$ set of nbhds arnd every pt $q$ forms an open cvr of $M$.
since $M$ cpt, then $\ex\set{U_{q_{1}},\dots,U_{q_{n}}}$ cvring $M$. consider:
\begin{equation*}
	G_{p}(x) \ = \ \max_{j}H_{pq_{j}}(x)
\end{equation*}
note $G_{p}\in\bar{A}$, and $G_{p}(x)>F(x)-\ep$ for all $x\in M$.
the first inequality has been satisfied.

now, $\fa p\in M$, $\ex$ nbhd $V_{p}\ni p$ s.t. $\fa x\in V_{p}$:
\begin{equation*}
	G_{p}(x)<F(x)+\ep
\end{equation*}
as before, $V_{p_{1}},\dots,V_{p_{m}}$ cvr $M$. let:
\begin{equation*}
	G \ = \ \min_{j}G_{p_{j}}\in\bar{A}
\end{equation*}
similar to before, we get that $G(x)<F(x)+\ep$ for all $M$.
OTOH, $G(x)>F(x)-\ep$ for all $M$, so we have found our desired function.
\qed

we need to patch a small hole near the start:
in particular, $\bar{A}$ may not contain any constant fn's.
\begin{lm}
	sps $I=[-M,M]$. $\ex$ seq of polys w/ no constant term s.t.
	$p_{n}\rightrightarrows\abs{x}$ in $I$.
\end{lm} \

\begin{pf}[source=Primary Source Material]
	by weierstrass approx. thm, $\ex q_{n}(x)=a_{n}+p_{n}(x)$ s.t.
	$q_{n}\rightrightarrows\abs{x}$. then:
	\begin{equation*}
		q_{n}(0)\sto\abs{0}=0 \ \implies \ a_{n}=q_{n}(0)\sto0 \trm{ as } n\sto\infty
	\end{equation*}
	since $a_{n}\sto0$, then $p_{n}=q_{n}-a_{n}\rightrightarrows\abs{x}$.
\end{pf}

r we rly doing banach fixed pt lol.

\begin{defn}
	we say $f$ is a \textbf{trigonometric polynomial} if:
	\begin{equation*}
		f(x)=a_{0}+\sum_{k=1}^{n}a_{k}\cos(kx)+\sum_{k=1}^{m}b_{k}\sin(kx)
	\end{equation*}
	as finite sums.
\end{defn} \

\begin{thm}
	let $f$ be $2\pi$-periodic cts.
	then $\ex p_{n}$ seq of trig polys s.t. $p_{n}\rightrightarrows f$.
\end{thm}
want to use stone-weierstrass, so want to restrict to $[0,2\pi]$.
however, 0 and $2\pi$ are not separated, and we can't exclude either pt while
maintaining cptness. idea: we identify them together instead.

\newpage
\begin{pf}[source=Primary Source Material]
	Identify the space of $2\pi$-periodic cts $f:\bR\gto\bR$ w cts
	$\tilde{f}:S^{1}\gto\bR$ by setting $\tilde{f}(e^{i\theta})=f(\theta)$ for $\theta\in[0,2\pi]$.

	need to check trig polys satisfy hypotheses as elems of $C_{0}(S^{1})$.
	indeed, they're an alg bc of trig identities, nowhere vanishing b/c they
	include constant fn's, and they separate pts (exercise).

	thus, by stone-weierstrass, we're done.
\end{pf}

hm. weird defn of a contraction. oh i lied its standard
\begin{defn}
	$f:M\gto M$ a \textbf{contraction} if $\ex\alpha\in[0,1)$ s.t.
	$d(f(x),f(y))\leq\alpha d(x,y)$ on all $M$.
\end{defn}

statement of banach fixed pt for completion
\begin{thm}[title=Banach Fixed-Point Theorem]
	sps $M$ complete, $f:M\gto M$ a contraction.
	then $\uex x\in M$ s.t. $f(x)=x$.
\end{thm} \

\begin{xmp}[source=Primary Source Material]
	let $(a_{ij})$ be doubly indexed real nums s.t.
	$\sup_{i,j\geq1}\abs{a_{ij}}\leq K$ and $\sum_{i,j}\abs{a_{ij}}<\infty$
	for some $K>0$. let $b\in\ell^{1}$ and consider:
	\begin{equation*}
		F:\ell^{1}\gto \ell^{1} \qquad (F(x))_{i}=\sum_{j}a_{ij}x_{j}^{2}+b_{j}
	\end{equation*}
	then:
	\begin{equation*}
		\sum_{i}(F(x))_{i}
		\ \leq \ \sum_{i}\left(\sum_{j}\abs{a_{ij}}\abs{x_{j}}^{2}+\abs{b_{i}}\right)
		\ \leq \ \norm{x}_{\infty}^{2}\sum_{ij}\abs{a_{ij}}+\abs{b}_{1}<\infty
	\end{equation*}
	we have $F:\ell^{1}\gto\ell^{1}$ is well-defn'd.
	we show that $\ex\ep>0$ s.t. if $\norm{b}_{1}\leq\ep$, then
	$\ex x\in\ell^{1}$ s.t. $x=F(x)$.

	idea: wtfind closed set $\subseteq\ell^{1}$ on which $F$ contracts.
	let $x,y\in\ell^{1}$. then:
	\begin{align*}
		\norm{F(x)-F(y)}_{1}
		& \ \leq \ \sum_{i}\abs{\sum_{j}a_{ij}(x_{j}^{2}-y_{j}^{2})} \\
		& \ \leq \ \sum_{i}\sum_{j}\abs{a_{ij}}\left(\abs{x_{j}}+\abs{y_{j}}\right)\abs{x_{j}-y_{j}} \\
		& \ \leq \ \norm{x-y}_{\infty}(\norm{x}_{\infty}+\norm{y}_{\infty})\left(\sum_{i,j}\abs{a_{ij}}\right) \\
		& \ \leq \ \norm{x-y}_{1}(\norm{x}_{1}+\norm{y}_{1})M
	\end{align*}
	note that as long as $x,y\in \bar{B}(0,r)$ with $r=1/(3M)$, then $F$ is a contraction.
	lastly, need to check that $F(\bar{B}(0,r))\subseteq\bar{B}(0,r)$.
	\begin{align*}
		\norm{F(x)}_{1}
		& \ \leq \ \sum_{i}\abs{\sum_{j}a_{ij}x^{2}+b_{i}} \\
		& \ \leq \ \norm{x}_{\infty}^{2}\sum_{i,j}\abs{a_{ij}}+\norm{b}_{1} \\
		& \ \leq \ \norm{x}_{1}^{2}M + \norm{b}_{1} \\
		& \ \leq \ \frac{M}{9M^{2}}+\ep
	\end{align*}
	whered that last bit come from. what.
	anyway as long as $M\geq1$ and $\ep\leq1/(9M)$:
	\begin{equation*}
		\norm{F(x)}_{1} \ \leq \ \frac{1}{9M}+\frac{1}{9M} \ \leq \ \frac{1}{3M}
	\end{equation*}
	so we can apply banach fixed pt.
\end{xmp}
point of the example: given non-linear fn's on a general space(?), it
may not be a contraction on the whole space.
however, we might be able to find a suff. small closed ball on which it is
in fact a contraction; if the fn also maps the ball to itself, then we're good.




